\section{Funcional aditivo e representação canônica}
\label{sec:functional}

\subsection{Dados racionais e perda local}

Considere observações
\begin{equation}
 \mathcal Z=\{(x_i,z_i,\omega_i):1\le i\le P\},
 \qquad
 x_i\in K\cap\mathbb Q,\quad z_i\in\mathbb Q,\quad
 \omega_i\in\mathbb Q_{>0}.
 \label{eq:rational-data}
\end{equation}
Para que uma observação em fronteira seja contada uma única vez, defina a
parte proprietária de $J=[a,b]\subseteq K=[\alpha,\beta]$ por
\begin{equation}
 J^{\own,K}=
 \begin{cases}
 [a,b),&b<\beta,\\
 [a,b],&b=\beta.
 \end{cases}
 \label{eq:relative-ownership}
\end{equation}
As partes proprietárias das folhas de qualquer árvore particionam $K$.
Defina a perda racional exata
\begin{equation}
 \loss_{\mathcal Z}(J;u,v)
 =\sum_{i:\,x_i\in J^{\own,K}}
 \omega_i\bigl(z_i-\ell_{J;u,v}(x_i)\bigr)^2.
 \label{eq:local-loss}
\end{equation}
Pela continuidade nas fronteiras, a soma dessas perdas sobre as folhas é a
perda quadrática empírica da montagem e independe da convenção de posse.

A teoria abaixo aceita qualquer família de custos locais não negativos em um
corpo ordenado computável. A escolha \eqref{eq:local-loss} é mantida porque
fecha uma implementação inteiramente racional. Erros integrais ou uniformes
podem ser usados quando um oráculo local exato ou certificado estiver
disponível; a \cref{sec:hard-budgets} trata explicitamente o caso uniforme.

\subsection{Fluxo recursivo}

Se $a$ é um átomo, denote por $n^\sharp(a)$ seu posto em $\Gamma^\sharp$ e
por $\gamma(q)$ o código gama de Elias do inteiro positivo $q$. Defina
recursivamente um fluxo binário para representações:
\begin{equation}
 \ser_{\rm rec}(R)=
 \begin{cases}
 1\,\gamma(n^\sharp(a)+1),&R\text{ é uma folha rotulada por }a,\\
 0\,\ser_{\rm rec}(R_L)\ser_{\rm rec}(R_R),
 &R\text{ tem as subárvores ordenadas }R_L,R_R.
 \end{cases}
 \label{eq:recursive-stream}
\end{equation}
Esse fluxo fornece ao seletor uma representação recursiva prefix-free. A
serialização condicionada da \cref{sec:complexity} permanece como o formato
usado nas cotas de compressão.

\begin{lemma}[decodificação e aditividade]
\label{lem:recursive-stream}
O conjunto dos fluxos \eqref{eq:recursive-stream} é prefix-free e cada fluxo
recupera uma única árvore rotulada. Seu comprimento $B(R)$ satisfaz
\begin{equation}
 B(a)=1+|\gamma(n^\sharp(a)+1)|,\qquad
 B(R_L\vee R_R)=1+B(R_L)+B(R_R).
 \label{eq:stream-additivity}
\end{equation}
\end{lemma}

\begin{proof}
O primeiro bit distingue folha de subdivisão. No primeiro caso, o código gama
é autodelimitado e recupera o endereço do único átomo. No segundo, o
decodificador chama a si mesmo duas vezes e consome exatamente as duas
subárvores ordenadas. A indução no número de nós prova unicidade da
decodificação. Se um fluxo válido fosse prefixo próprio de outro, o
decodificador terminaria antes do fim do segundo fluxo, contradizendo o
consumo exato de uma árvore. A contagem dos bits dá
\eqref{eq:stream-additivity}.
\end{proof}

\subsection{Objetivo e desempate}

Fixe penalidades racionais
\[
 \lambda_{\rm bits}\ge0,\qquad \lambda_{\rm split}\ge0.
\]
Para uma representação $R$, com função avaliada $g_R$, defina
\begin{equation}
 \cost_{\mathcal Z}(R)
 =\sum_{J\in\Lambda(R)}
   \loss_{\mathcal Z}\bigl(J;g_R(\inf J),g_R(\sup J)\bigr)
 +\lambda_{\rm bits}B(R)
 +\lambda_{\rm split}\,I(R),
 \label{eq:additive-objective}
\end{equation}
onde $\Lambda(R)$ é o conjunto de folhas e $I(R)$ é o número de nós internos.
Em uma árvore binária cheia, $I(R)=|\Lambda(R)|-1$, mas manter a notação
separada deixa visível o custo de uma decisão de subdivisão.

Perda quadrática aditiva e penalização da complexidade sobre partições
recursivas aparecem em CART e em árvores diádicas ótimas
\citep{breiman1984,donoho1997,blanchard2007}. Os termos de comprimento binário,
a admissibilidade aritmética das folhas e o desempate canônico em
\eqref{eq:additive-objective} são específicos do problema presente.

Para um par admissível $(u,v)\in\Pairs_H(J)$, seja $a(J;u,v)$ o único átomo
da \cref{prop:endpoint-dictionary}. Defina o custo de folha
\begin{equation}
 \Phi_H(J;u,v)=
 \loss_{\mathcal Z}(J;u,v)
 +\lambda_{\rm bits}
  \bigl(1+|\gamma(n^\sharp(a(J;u,v))+1)|\bigr),
 \label{eq:leaf-cost}
\end{equation}
e ponha $\Phi_H(J;u,v)=+\infty$ quando $(u,v)\notin\Pairs_H(J)$.

Empates são resolvidos pela chave total
\begin{equation}
 \kappa(R)=
 \bigl(\cost_{\mathcal Z}(R),\,
       |\ser_{\rm rec}(R)|,\,
       \ser_{\rm rec}(R)\bigr),
 \label{eq:tie-key}
\end{equation}
ordenada lexicograficamente, usando a ordem lexicográfica usual no último
componente. Assim, primeiro minimiza-se o funcional; depois, o fluxo mais
curto; por fim, o menor fluxo binário.

\begin{lemma}[compatibilidade do desempate com a colagem]
\label{lem:tie-concatenation}
Considere subárvores compatíveis com os mesmos estados externos, de modo que
as substituições abaixo preservem a continuidade. Se
$\kappa(R_L)\le\kappa(R'_L)$ e $\kappa(R_R)\le\kappa(R'_R)$, então
\[
 \kappa(R_L\vee R_R)\le\kappa(R'_L\vee R'_R).
\]
Em particular, substituir um filho por outro de chave menor nunca piora a
chave de nenhuma superárvore.
\end{lemma}

\begin{proof}
Na colagem, o custo é a soma dos custos dos filhos mais a mesma constante
$\lambda_{\rm bits}+\lambda_{\rm split}$, e o comprimento é
$1+B(R_L)+B(R_R)$. Logo uma desigualdade estrita em custo, ou depois em
comprimento, é preservada. Se custo e comprimento empatam, a comparação final
é entre
$0\,\ser_{\rm rec}(R_L)\ser_{\rm rec}(R_R)$ e
$0\,\ser_{\rm rec}(R'_L)\ser_{\rm rec}(R'_R)$. Como fluxos que chegam a essa
terceira coordenada têm os mesmos comprimentos nos blocos substituídos, a
ordem lexicográfica é decidida primeiro no bloco esquerdo que difere e depois
no direito. Portanto ela também é preservada. Aplicar o argumento a um filho
de cada vez prova a afirmação.
\end{proof}
